\documentclass[10pt,a4paper]{amsart}
\usepackage[margin=19mm]{geometry}
\usepackage{amsmath,amssymb,mathtools}
\usepackage[scr=rsfs,cal=euler]{mathalfa}
\usepackage{xfrac}
\usepackage[unicode,hidelinks,bookmarks=false]{hyperref}
\newcommand{\ie}{i.e.\ }
\newcommand{\cf}{cf.\ }
\newcommand{\resp}{respectively\ }
\newcommand{\Subsection}{\subsection}
\providecommand{\nobreakdash}{\nobreak\hbox{-}}
\setlength{\emergencystretch}{2em}
\multlinegap0pt
\usepackage{amssymb}
\newtheorem{thm}{Theorem}
\newtheorem{lem}{Lemma}
\newcommand{\Om}{\Omega}
\newcommand{\bdry}{\partial}
\newcommand{\bdy}{\bdry}
\DeclareMathOperator*{\essliminf}{ess\,lim\,inf}
\newcommand{\p}{{$p\mspace{1mu}$}}
\title{Liouville theorems and removable sets for bounded $p$-harmonic and quasiharmonic functions on metric spaces under local assumptions}
\author{Anders Bj\"orn, Jana Bj\"orn}
\date{Source: arXiv:2608.26878v1 (CC BY 4.0)}
\begin{document}
\maketitle

\noindent\emph{Source and licence.} Liouville theorems and removable sets for bounded $p$-harmonic and quasiharmonic functions on metric spaces under local assumptions. Version arXiv:2608.26878v1 (CC BY 4.0), distributed by arXiv under the Creative Commons Attribution 4.0 International licence, http://creativecommons.org/licenses/by/4.0/, as recorded in the arXiv licence metadata for this exact version and shown on the abstract page https://arxiv.org/abs/2608.26878v1. Full licence notice: \texttt{licenses/arxiv-2608.26878-CC-BY-4.0.txt}.\par\smallskip

\noindent\textit{Editorial scope.} Counted unit: the article's thm thm-Liouville-para (source0.tex lines 842--851). The article's complete original proof of it is delivered with it (source0.tex lines 2013--2025, one \texttt{proof} environment of 13 lines, which the article places away from the statement). No mathematics is written, repaired or re-derived: the statement and the proof are the article's own lines, in the article's own notation, and no retained line is substituted or re-flowed.  Retained uncounted as the source-local context the proof needs: lines 1856--1858; lines 1917--1933.  Nothing was omitted from any retained span, so the package carries no omission record.  No retained line contains a citation, so the wrapper carries no bibliography and no bibliography entry.  No retained line contains a figure environment, an image-inclusion command or drawing code, so no image is delivered and the package declares no asset dependency.  Editorial formatting only, in addition: an AMS \texttt{amsart} preamble that restates the article's own declarations and macros used by the retained lines, namely the environments \texttt{thm}, \texttt{lem} on separate counters, together with the standard packages those lines need.  The retained environments print in this document as Theorem 1, Lemma 1 in order of appearance, whereas the article prints the same items with the numbering of its own full document.  The complete native source of the article as distributed in the arXiv e-print for this exact version is bundled unchanged as \texttt{sources/208670/source0.tex}.
\begin{equation} \label{eq-essliminf-reg}
u(x)=\essliminf_{y \to x} u(y) \quad \text{for } x \in \Om,
\end{equation}

\begin{lem}   \label{lem-quasi-Liouville-Om}
Assume that $\bdy\Om = \{x_0\}$ and that $\Om$ is either
bounded or a \p-parabolic set.
Let $u$ be a 
quasisuperharmonic function in $\Om$ that is bounded from below,
and let
\begin{equation}   \label{eq-def-u-x0}
u(x_0)=\liminf_{\Om\ni y\to x_0} u(y).
\end{equation}
Then $u(x_0)<\infty$ and $u\ge u(x_0)$ in $\Om$.

In particular, if $u$ is a bounded quasiharmonic function in $\Om$, such that 
\begin{equation}   \label{eq-lim-u-x0}
\lim_{\Om\ni y\to x_0} u(y)\quad \text{exists},
\end{equation}
then $u$ is constant.
\end{lem}

\begin{thm} \label{thm-Liouville-para}
\textup{(The Liouville theorem in \p-parabolic spaces)}
Assume that   $X$ is a proper connected metric space
equipped with a locally doubling measure $\mu$ that
supports a local \p-Poincar\'e inequality, where $1<p<\infty$.

If $X$ is \p-parabolic or bounded, then 
every nonnegative quasisuperharmonic function in $X$
is constant.
\end{thm}

\begin{proof}[Proof of Theorem~\ref{thm-Liouville-para}]
Let $\{x_j\}_{j=1}^\infty$ be a sequence such that 
\begin{equation} \label{eq-L-para} 
u(x_j)\to \sup_X u,
\quad \text{as } j\to\infty.
\end{equation}
Since $u$ is $\essliminf$-regularized
as in~\eqref{eq-essliminf-reg}, we see that
$u(x_j)=\liminf_{\Om\ni y\to x_j} u(y)$.
Lemma~\ref{lem-quasi-Liouville-Om} implies that
$u  \ge u(x_j)$ in $X$ for $j=1,2,\dots$\,.
Together with~\eqref{eq-L-para} this shows that $u \equiv \sup_X u$.
\end{proof}

\end{document}
